\documentclass[11pt,a4paper]{article}
\usepackage[utf8]{inputenc}
\usepackage[T1]{fontenc}
\usepackage{lmodern}
\usepackage[brazil]{babel}
\usepackage{amsmath,amssymb,amsthm}
\usepackage{geometry}
\usepackage{hyperref}
\usepackage{booktabs}
\usepackage{enumitem}
\usepackage{xcolor}

\geometry{margin=2.5cm}

\definecolor{remark}{rgb}{0.12,0.37,0.55}

\newtheorem{theorem}{Proposição}
\newcommand{\E}{\mathbb{E}}
\newcommand{\dd}{\mathrm{d}}

\title{\textbf{Capítulo 8 --- Consumo}\\[4pt]
       {\large Renda Permanente, Passeio Aleatório e Poupança Precaucional}\\[2pt]
       {\normalsize PPGEco-CCJE-UFES · Macroeconomia I · Romer (2019) Cap.\ 8}}
\author{Projeto de Estudo --- \textit{macro-romer}}
\date{2026}

\begin{document}
\maketitle
\tableofcontents
\bigskip

% ============================================================
\section{O problema intertemporal do consumidor}
% ============================================================

Um consumidor escolhe $\{c_t\}_{t=0}^{\infty}$ para maximizar a utilidade
esperada descontada
\begin{equation}\label{eq:objective}
  \E_0 \sum_{t=0}^{\infty} \beta^t\, u(c_t), \qquad \beta = \frac{1}{1+\rho},
\end{equation}
sujeito à restrição orçamentária dinâmica
\begin{equation}\label{eq:budget}
  a_{t+1} = (1+r)(a_t + y_t - c_t),
\end{equation}
onde $a_t$ são ativos no início do período, $y_t$ é a renda (do trabalho)
e $r$ é a taxa real de juros (constante).

\paragraph{Restrição orçamentária intertemporal.}
Iterando \eqref{eq:budget} para a frente e impondo a condição de
não-Ponzi $\lim_{T\to\infty} a_T/(1+r)^T = 0$, obtém-se
\begin{equation}\label{eq:lifetime_budget}
  \sum_{t=0}^{\infty} \frac{c_t}{(1+r)^t}
  = (1+r)a_0 + \sum_{t=0}^{\infty} \frac{y_t}{(1+r)^t}
  \;\equiv\; (1+r)a_0 + h_0,
\end{equation}
onde $h_0$ é a \emph{riqueza humana} --- o valor presente do fluxo de
renda do trabalho.

% ============================================================
\section{Hipótese da Renda Permanente (PIH)}
% ============================================================

\subsection{Caso determinístico: consumo como anuidade}

Suponha $\beta(1+r) = 1$ e utilidade não necessariamente quadrática, mas
com renda perfeitamente prevista $\{y_t\}$. A condição de primeira ordem
(Euler) é
\[
  u'(c_t) = \beta(1+r)\,u'(c_{t+1}) = u'(c_{t+1})
  \;\Longrightarrow\; c_t = c_{t+1} = c \;\;\forall t.
\]
Substituindo $c_t \equiv c$ em \eqref{eq:lifetime_budget} e usando
$\sum_t (1+r)^{-t} = (1+r)/r$:
\begin{equation}\label{eq:pih_consumption}
  \boxed{c = \frac{r}{1+r}\Bigl[(1+r)a_0 + h_0\Bigr]
         = r\,a_0 + \frac{r}{1+r}\,h_0.}
\end{equation}

O consumo é a \textbf{anuidade} da riqueza total (financeira + humana):
o agente ``consome os juros'' da sua riqueza completa. Choques
\emph{transitórios} a $y_t$ alteram $h_0$ pouco (são divididos pelo
fator $(1+r)^t$ e, além disso, por toda a vida útil), de modo que $c$
responde pouco --- esta é a \textbf{suavização do consumo}.

\subsection{Implementação: \texttt{PermanentIncomeModel}}

No código (\texttt{ch08\_consumption.py}), a renda segue um AR(1) em
torno da média $\bar y$:
\[
  y_t - \bar y = \rho_y\,(y_{t-1} - \bar y) + \varepsilon_t,
  \qquad \varepsilon_t \sim N(0,\sigma_y^2).
\]
A riqueza humana de um desvio AR(1) tem solução fechada:
\[
  h_t = \frac{(1+r)\bar y}{r}
        + \frac{y_t - \bar y}{1 - \dfrac{\rho_y}{1+r}},
\]
pois $\E_t[y_{t+j} - \bar y] = \rho_y^j (y_t - \bar y)$ e
$\sum_{j\ge 0} \bigl(\rho_y/(1+r)\bigr)^j = \bigl(1-\rho_y/(1+r)\bigr)^{-1}$.
O consumo segue então de \eqref{eq:pih_consumption} com $h_0 \to h_t$.
A função \texttt{simulate} acumula $a_{t+1}=(1+r)(a_t+y_t-c_t)$ e
\texttt{random\_walk\_test} verifica que $\Delta c_t$ não é
autocorrelacionado.

% ============================================================
\section{Hall (1978): o consumo como passeio aleatório}
% ============================================================

\subsection{Utilidade quadrática}

Hall (1978) assume utilidade quadrática
\[
  u(c) = c - \frac{b}{2}c^2, \qquad u'(c) = 1 - bc,
\]
linear na utilidade marginal. A equação de Euler é
\[
  u'(c_t) = \beta(1+r)\,\E_t[u'(c_{t+1})].
\]
Sob a hipótese central $\beta(1+r) = 1$:
\[
  1 - bc_t = \E_t[1 - bc_{t+1}] = 1 - b\,\E_t[c_{t+1}]
  \;\Longrightarrow\;
  \boxed{c_t = \E_t[c_{t+1}].}
\]

\subsection{Consumo é uma martingala}

Definindo $c_{t+1} = c_t + \varepsilon_{t+1}$ com
$\E_t[\varepsilon_{t+1}] = 0$, o resultado de Hall implica que
\textbf{nenhuma informação disponível em $t$} --- exceto $c_t$ ---
ajuda a prever $c_{t+1}$. Em particular, variáveis defasadas como
$y_{t-1}, a_{t-1}$ não devem ter poder preditivo sobre $\Delta c_t$
condicional em $c_{t-1}$.

\paragraph{Implicações testáveis:}
\begin{itemize}
  \item $\mathrm{Cov}(\Delta c_t, \Delta c_{t-1}) = 0$ (sem autocorrelação).
  \item $\mathrm{Cov}(\Delta c_t, I_{t-1}) = 0$ para qualquer informação
    $I_{t-1}$ disponível em $t-1$.
  \item Variações de política fiscal \emph{antecipadas} não devem afetar
    $\Delta c_t$ no momento em que ocorrem (apenas no anúncio).
\end{itemize}

No código, \texttt{HallRandomWalk.test\_martingale\_property} faz Monte
Carlo de $\E[\Delta c_t]\approx 0$, e \texttt{euler\_holds} verifica
$\beta(1+r)=1$.

% ============================================================
\section{Sensibilidade excessiva e o modelo de Campbell-Mankiw}
% ============================================================

\subsection{A evidência de ``sensibilidade excessiva''}

Diversos estudos empíricos (Flavin 1981; Campbell e Mankiw 1989) rejeitam
a previsão de Hall: $\Delta c_t$ \textbf{é} correlacionado com
$\Delta y_t$ contemporâneo (e até com $y_{t-1}$). O consumo ``reage demais''
a mudanças de renda corrente.

\subsection{O modelo de duas populações}

Campbell e Mankiw (1989) propõem que uma fração $\lambda \in [0,1]$ da
população consome simplesmente sua renda corrente (``rule of thumb'',
sem acesso a crédito ou sem otimização intertemporal), enquanto a fração
$1-\lambda$ segue Hall:
\[
  c_t^{\text{RT}} = y_t
  \qquad\Longrightarrow\qquad \Delta c_t^{\text{RT}} = \Delta y_t,
\]
\[
  c_t^{\text{PIH}} \;\Longrightarrow\; \Delta c_t^{\text{PIH}} = \varepsilon_t,
  \quad \E_{t-1}[\varepsilon_t] = 0.
\]
Agregando,
\begin{equation}\label{eq:campbell_mankiw}
  \boxed{\Delta c_t = \lambda\,\Delta y_t + (1-\lambda)\,\varepsilon_t.}
\end{equation}

\paragraph{Interpretação de $\lambda$:}
\begin{itemize}
  \item $\lambda = 0$: PIH/Hall pura --- nenhuma sensibilidade excessiva.
  \item $\lambda = 1$: toda a população é ``mão-à-boca'' (hand-to-mouth);
    o consumo agregado replica a renda agregada.
  \item Estimativas empíricas tipicamente situam $\lambda \in [0{,}2, 0{,}5]$
    para economias desenvolvidas, e valores mais altos em economias com
    mercados de crédito menos profundos (como o Brasil).
\end{itemize}

A regressão \eqref{eq:campbell_mankiw} é estimada por OLS em
\texttt{CampbellMankiw.estimate\_lambda} e aplicada a dados brasileiros
(consumo das famílias e PIB, IBGE/SIDRA tabela 1846) em
\texttt{ch08\_consumption\_empirics.py}.

% ============================================================
\section{Poupança precaucional e o modelo buffer-stock}
% ============================================================

\subsection{Prudência: $u'''(c) > 0$}

Com incerteza sobre a renda futura, a equação de Euler geral é
\[
  u'(c_t) = \beta(1+r)\,\E_t[u'(c_{t+1})].
\]
Por Jensen, se $u'$ é \textbf{convexa} ($u'''>0$, ``prudência'' no sentido
de Kimball 1990), então
\[
  \E_t[u'(c_{t+1})] > u'(\E_t[c_{t+1}]),
\]
de modo que, para satisfazer a Euler, o consumidor escolhe $c_t$ menor
(poupa mais) do que escolheria sob certeza. Esse é o motivo
\textbf{precaucional} para poupar: a incerteza por si só eleva a
poupança, mesmo mantendo $\E_t[y_{t+1}]$ constante. A utilidade CRRA
$u(c)=c^{1-\theta}/(1-\theta)$ tem $u'''(c) = \theta(\theta+1)c^{-\theta-2}>0$
para $\theta>0$, portanto é prudente.

\subsection{Restrição de crédito e o problema com VFI}

O modelo \texttt{BufferStockModel} resolve, por iteração da função-valor
em uma grade de ativos $a \in [0, a_{\max}]$ e dois estados de renda
$y \in \{y_{\text{baixo}}, y_{\text{alto}}\}$ (cadeia de Markov com
persistência $\pi$):
\begin{equation}\label{eq:bellman}
  V(a,y) = \max_{0 \,\le\, a' \,\le\, m(a,y)}
    \Bigl\{u\bigl(m(a,y) - a'\bigr) + \beta\,\E\bigl[V(a',y')\mid y\bigr]\Bigr\},
\end{equation}
onde $m(a,y) = (1+r)a + y$ é a ``liquidez na mão'' (cash-on-hand) e a
restrição $a' \ge 0$ impede empréstimos (sem crédito ao consumidor).

\paragraph{Condição de impaciência.}
Para que exista um ``alvo'' de riqueza estacionário (e o agente não
acumule ativos indefinidamente), é necessário $\beta(1+r) < 1$: o agente
é suficientemente impaciente para que, na ausência de incerteza, ele
desejaria \emph{des}-acumular ativos. A incerteza de renda, porém, gera
um motivo precaucional que compensa essa impaciência perto de $a=0$,
produzindo um \textbf{estoque-tampão} (buffer stock) de ativos ao redor
do qual a riqueza flutua estacionariamente.

\subsection{Função-política}

A solução numérica gera $c^*(a,y)$, que satisfaz:
\begin{itemize}
  \item $c^*(a,y) \le m(a,y)$ (consumo nunca excede os recursos);
  \item $c^*(a,y)$ é crescente em $a$ e em $y$;
  \item para $a$ próximo de $0$, $c^*(a,y) < m(a,y)$ estritamente
    (poupança precaucional positiva mesmo perto da restrição).
\end{itemize}

A simulação de painel (\texttt{simulate\_panel}) gera uma distribuição
ergódica de ativos entre famílias, ilustrando a heterogeneidade gerada
puramente por choques idiossincráticos de renda.

% ============================================================
\section{Síntese}
% ============================================================

\begin{table}[h!]
  \centering
  \caption{Modelos de consumo --- resumo}
  \begin{tabular}{lll}
    \toprule
    Modelo & Hipótese chave & Implicação testável \\
    \midrule
    PIH (Friedman) & $\beta(1+r)=1$, certeza equiv. &
      $c$ = anuidade da riqueza total; choques transitórios $\to$ $c$ quase fixo \\
    Hall (1978) & Utilidade quadrática + $\beta(1+r)=1$ &
      $c_t = \E_t[c_{t+1}]$: $\Delta c_t$ imprevisível \\
    Campbell-Mankiw & Fração $\lambda$ ``mão-à-boca'' &
      $\Delta c_t = \lambda \Delta y_t + (1-\lambda)\varepsilon_t$, $\lambda>0$ \\
    Buffer-stock & $\beta(1+r)<1$ + $u'''>0$ + $a'\ge 0$ &
      Poupança precaucional; distribuição ergódica de riqueza \\
    \bottomrule
  \end{tabular}
\end{table}

\noindent A trajetória teórica do capítulo --- de Friedman a Hall, e
de Hall às suas rejeições empíricas (Campbell-Mankiw, buffer-stock) ---
espelha a trajetória de Fisher a Calvo no Capítulo~7: um modelo de
referência elegante (PIH/Hall, Calvo-básico) é progressivamente
enriquecido com fricções realistas (restrição de crédito, prudência,
heterogeneidade) que melhoram o ajuste aos dados sem descartar a
intuição original.

\end{document}
